% =====================================================================
%  Ringraziamenti -- capitolo non numerato, incluso in coda alla tesi
%  (dopo la bibliografia) da anteprima_bozze_it.tex via \input.
% =====================================================================
\cleardoublepage
\phantomsection
\addcontentsline{toc}{chapter}{Ringraziamenti}
\chapter*{Ringraziamenti}
\markboth{Ringraziamenti}{Ringraziamenti}

Eccoci qui arrivati al tanto atteso capitolo dei ringraziamenti.
Chissà come verrà percepito: chi piangerà nel sentirli, chi ci farà una
risata. Oggi è un giorno speciale per il sottoscritto, segna una delle
pagine più importanti della mia vita: la prima laurea; un titolo che
rappresenta prestigio, sacrificio, dedizione e virtù.

\section{Parenti}
Ora è il momento di fare qualche passo indietro e dare la giusta
dignità e merito a chi mi ha sostenuto fino ad oggi. Tra tutti, i primi
sono senz'altro i genitori, che sono il primo punto di riferimento --
il baluardo -- di qualsiasi persona. Ci sono da prima ancora che ci si
possa effettivamente dichiarare ``nati'', perché veniamo generati
proprio a partire da loro: contengono entrambi parte del codice
sorgente della nostra vita, che viene poi fuso insieme, risolto con
tutte le dipendenze legate alle funzioni biologiche e infine, dopo 9
mesi, durante i quali viene compiuta un'attenta analisi di debug
all'ospedale, viene rilasciata la build finale funzionante: ovvero noi.
Già solo per questo meritano un enorme grazie, un grazie per avermi
committato e pushato su questo mondo. Oltre a quanto detto -- in
maniera molto allegorica --, sono a loro molto riconoscente per avermi
dato l'opportunità di crescere bene e per avermi fatto sperimentare, in
tale contesto, le mie passioni senza mai farmi mancare nulla.

\subsection{Pap\`a}
Ringrazio papà per avermi messo in grado di riconoscere alcune delle
cose essenziali della vita: gli affetti, le regole, la gratitudine e la
forza di andare avanti, oltre le mille mila difficoltà, che il più
delle volte siamo noi stessi a ingigantire troppo. Con lui ho discusso
spesso di come alle volte mi sentivo schiacciato dalle paranoie sul
futuro ancora lontano, di come queste cose mi divoravano e non mi
facevano stare sereno; lui me le ha fortemente ridimensionate,
ricordandomi che nulla può essere certo. Non sappiamo come operano le
tre Moire lassù, con in mano il nostro fragile filo. Abbiamo potere
solo sulle cose presenti che possiamo controllare veramente. Lui ha
vissuto momenti veramente duri, talmente duri che non ho affatto
certezza che qualcun altro al posto suo potesse uscirne a testa alta
come ha fatto lui.

Papà, quando era giovane, non ha avuto le stesse opportunità che ho
avuto io; dovette fare molti sacrifici per arrivare a soddisfare i suoi
desideri, senza permettersi mai di lamentarsi. Non aveva sicurezze
circa il futuro, ma sapeva che l'unica cosa che doveva continuare a
fare era vivere seguendo i suoi principi, mettendo le responsabilità al
primo posto.

Lui ha vissuto facendo della fatica il suo motore catalizzatore,
impegnandosi nel suo duro lavoro con serietà, professionalità e
talento. Oggi nel suo settore eccelle più di chiunque altro perché è
stato scaltro: ha reso ciò che faceva la sua passione e per questo non
gli è mai mancata la voglia di fare. Dopo la sua lunga carriera ancora
spinge come un treno e ha più energia in corpo di quanta se ne possa
vedere in ragazzi di 30 anni più giovani. Lui, al contrario di me, vive
la vita con più leggerezza: non rinuncia a commettere sgarri nella
dieta abituale, non fa caso a tutti quei micro dettagli che a me
manderebbero fuori di testa e non si fa problemi nel vivere solo il
momento presente; dimostra tra l'altro un fisico migliore del mio e un
carattere più mite e sereno.

Abbiamo molte differenze sostanziali, opinioni differenti su tante
cose, ma nonostante ciò lo venero comunque. Rappresenta per me la
manifestazione dei valori e delle virtù che dovrebbe avere un uomo:
essere dedito alle responsabilità; avere pazienza, in quanto essa è la
virtù dei forti -- come è solito dirmi sempre --; stare bene con sé
stesso e i propri cari; non essere mai mediocre in ciò che si sceglie
di fare.

\paragraph{Menzione d'onore}
Fu lui a regalarmi la mia prima console da gioco: il
leggendario Nintendo DS. Da allora nacque la mia grande passione per il
medium videoludico, che mi accompagna fino a oggi. Grazie ancora pa'.

\subsection{Mamma}
Adesso passiamo alla mamma. Lei che più di tutti ci teneva a questo mio
traguardo. Ci vorrebbero tante altre pagine -- quasi a raggiungere il
doppio di quelle dedicate alla tesi -- per raccontare tutti i suoi
meriti, ma cerco di stringermi per non tediare troppo anche gli altri,
che scommetto staranno morendo dal caldo e vorrebbero in cuor loro che
io accelerassi di un ordine di grandezza.

Ringrazio mamma in primis per l'educazione e la disciplina che mi ha
trasmesso. Quando ero piccolo non ero affatto contento di finire in
punizione quasi ogni settimana, a causa delle note di merito e dei
brutti voti che portavo a casa -- sin dalle elementari. Mi dava grande
fastidio la sua severità: rispetto alle altre mamme dei miei amici,
finiva sempre per sgridarmi e punirmi anche per le cose apparentemente
più infime e sciocche. Ad oggi dico grazie, un grazie grande come una
nazione. Sono grato per il fatto che mi veniva tolta la play quando non
avevo voglia di studiare o esageravo con la durata delle sessioni; sono
grato per il cibo spazzatura che mi veniva vietato quando iniziavo a
mettere su qualche chilo di troppo; sono grato per gli schiaffi che
ricevevo se rispondevo male o se mi permettevo di essere maleducato con
il prossimo; sono grato per le punizioni che mi venivano impartite per 
un brutto voto di una verifica.

Su certe cose tanti ci sarebbero passati sopra, ma la mia mamma non
transigeva su nulla. Nonostante le grandi difficoltà che le arrecava la
vita in certi momenti, ha sempre mantenuto la riga dritta nei miei
confronti.

Mi ha insegnato a sapermi comportare; mi ha spinto alla lettura,
ritenendo questa fondamentale perché io potessi attivare la testa; mi
ha fornito il metodo di studio pagando le ripetizioni private che
impegnavano gran parte delle settimane dei miei anni di elementari -- a
detta sua, con il totale complessivo di queste spese sarebbe riuscita a
comprarci un'utilitaria --; mi teneva alla larga da quelle
frequentazioni che lei riteneva di cattiva influenza. Chiaramente, come
in qualsiasi rapporto madre-figlio, su queste sue decisioni io non ero
mai d'accordo. Mi scocciava non poter più giocare o guardare i miei
videogiochi e cartoni preferiti ogni qual volta riportavo a casa
un'insufficienza. Non volevo leggere i libri che mi comprava, preferivo
di gran lunga guardare YouTube. Volevo che mi si comprassero i
pacchettoni di patatine da mangiare nel mentre che mi guardavo i film.
Non volevo partecipare ai campus estivi perché provavo un astio enorme
verso qualunque attività fisica.

Mia madre se ne è sempre infischiata delle mie ragioni e ha sempre
esercitato la sua autorità da genitore nel modo in cui riteneva più
giusto, facendolo quindi per il mio bene. Ad oggi, quali meriti posso
attribuirle dopo tutto ciò? Posso dire: di sapere come pormi con il
prossimo, mostrando rispetto e inclusione -- salvo per chi non ritengo
degno --; di mantenere una disciplina ferrea quando serve; di avere
una buona capacità linguistica e di conoscere abbastanza cose grazie
alle vagonate di libri che mi sono letto in questi anni, i quali non
avrei mai neanche sfogliato senza la sua spinta iniziale; di avere un
bel fisico dopo i tanti allenamenti di nuoto e di palestra promossi e
finanziati da lei. 

Così come papà, porta avanti le sue responsabilità, si dedica
al lavoro con passione ed è fortissima nel suo ambito, e si va persino
ad allenare in palestra anche dopo le giornate più faticose. Insieme
hanno lavorato tanto per forgiare il mio carattere e per aiutarmi nel
raggiungimento dei miei traguardi. È grazie a loro che oggi arrivo qui
a conseguire il titolo di laurea triennale con il massimo dei voti.
Spero vivamente di poter regalare loro ancora molte altre
soddisfazioni.

\paragraph{Menzione d'onore}
Ringrazio mamma anche per aver seguito una dieta
ricca di proteine durante la fase di building della mia immagine,
cosicché io abbia potuto vantare tutti i capelli in testa sin dal
giorno del mio concepimento.

\subsection{Nonni, Zii \& Cugini}
Non posso affatto esimermi dal menzionare nei ringraziamenti anche i
nonni, gli zii e i cugini che hanno orbitato attorno alla mia
esistenza, lasciandomi ognuno qualcosa di buono dentro. In particolare
la mia nonna materna, che mi ha visto letteralmente crescere giorno
dopo giorno. In tutti questi anni, dopo l'uscita da scuola o
dall'università, ho avuto l'appuntamento fisso a casa sua per l'ora di
pranzo. Il nostro rituale è stato anche uno dei forti motivi per i
quali ho deciso di continuare a fare il pendolare per seguire le
lezioni a Roma, così da poter prendere il primo treno utile per poterle
andare a fare compagnia. Per lei sono la cosa più preziosa che ha, il
suo unico nipote, quindi ci tengo a volerle dare il giusto rilievo in
questa sezione. Ti ringrazio, Franca, per tutti questi anni in cui mi
hai ``accudito'' e per il tuo eterno tifo; che possa il nostro rituale
durare ancora negli anni a venire.

Un abbraccio figurato è rivolto anche ai miei nonni paterni, ai quali
sono andato a fare visita ogni fine settimana insieme a papà.
Immancabilmente anche da loro sono stato sempre viziato con i pranzi
ricchi di varie portate. Marino e Giacomina hanno sempre mostrato il
loro interesse verso i miei traguardi e sono anche loro molto fieri di
me. Li ringrazio anche per il fatto di riuscire a soddisfare il mio
fabbisogno di una trentina di uova a settimana grazie al loro prolifico
pollaio; senza queste ultime non avrei avuto questo boost in
performance.

Ringrazio zia Alessandra, zio Mario e i miei cugini Andrea e
Nico, che condividono con me la gioia di questo giorno. Confermo
finalmente quello che mi diceva Andrea: oggi sono stato promosso a
``primo dottore in famiglia Pierbattisti''. Un abbraccio sentito a
tutti voi.

\section{Amici}
Adesso viene il turno degli amici. Gli amici sono le persone più
importanti della nostra vita dopo i genitori: ci plasmano, ci
sostengono, ci aiutano a passare meglio il tempo su questa Terra. È
difficile per me definire il concetto vero di amicizia, ma a
prescindere da ciò sono sicuro di poter affermare di averne molti su
cui contare oggi. Mi rendo conto che per tanti di loro io somigli a un
marziano, per via delle abitudini e convinzioni che ho e delle azioni
fuori dalla norma che compio; le nostre idee spesso divergono, ma ciò
che rimane importante è che troviamo il modo di stare bene insieme a
prescindere dal resto. Con il passare del tempo ho imparato ad
accettare le differenze fra noi persone, a saperci convivere e,
soprattutto, a tirare fuori il buono da ciascuno. Ormai ho solidificato
i rapporti con ognuno attraverso le bellissime esperienze
vissute insieme e la vostra presenza sincera nei miei riguardi. Posso
quindi ritenermi fortunato; ho amici che mi sostengono e mi migliorano,
sia sul piano dell'essere che dell'avere.

\subsection{Minecraftiani}
Ringrazio il trio Minecraftiani, il leggendario
gruppo creatosi da inizio liceo e che ancora oggi tiene unito il legame
tra i suoi membri: me medesimo (Pierba), Matteo (Dexter) e Riccardo
(Uffre). Il nome è dato dal fatto che questo gruppo si originò per il
nostro bisogno di giocare insieme all'interno di un server di
Minecraft. Di quel mondo condiviso conservo ancora tante bellissime
memorie: le serate di Discord, il martedì in cui Uffre sbloccava il
gaming diurno, le immense farm e costruzioni a cui ci dedicavamo come
fosse un lavoro. Nel tempo il nostro gruppo si evolse, cambiò la sua
identità, ma il nome rimase immutato sin dal 16 settembre del 2019.
Siamo rimasti legati per tutti gli anni di liceo e oltre; siamo
sopravvissuti a verifiche, interrogazioni, alla riproduzione del
Battistero di Firenze. Durante quegli anni il nostro gruppo era incluso
in uno più grande -- Stanza 502 --, ma il nostro trio è l'unico che è
uscito indenne dalla prova del tempo.

Eccoci che in uno schiocco di dita siamo passati dal preoccuparci di
come organizzare la base del nostro server a come sistemare il nostro
curriculum vitae. Assurdo come debba passare il tempo. Vi ringrazio per
tutto ciò che siete stati per me fino a oggi. Non dimenticherò mai i
giorni delle nostre biciclettate, dei just-eat, delle giuserpate, del
campeggio -- nel quale avete dimostrato grosse palle nel sopportarmi --
e di tanti altri eventi canonici che non sto qui a raccontare se non
voglio che questo file pesi un Terabyte.

Mi dispiace se in certe occasioni non sono stato presente, come spiego
a tutti: non ho più l'età per fare certe cose; comunque impegniamoci a
mantenere viva la fiamma del nostro gruppo. Vi voglio bene.

\subsection{Membri del Partito}
Ringrazio i Membri del Partito. Ormai questo
gruppo è divenuto il gruppo dei ``soliti''. Ci siamo riuniti nel '23 e
da allora abbiamo costruito nel tempo la nostra piccola famigliola.
Anche se ci sono stati conflitti, ci siamo sempre riconciliati e, anche
se la nostra capacità di organizzazione è fortemente da rivedere, il
nostro lo reputo un gruppo solido e sano, e sono felice di farne parte.
Grazie ai Lorenzi (uno grosso, l'altro secco), a Giulia tanto buona e
cara, a Gianni che allena pure gambe e alla sua ragazza Irene che dà
sempre carta bianca, a vossignoria Buso, a papà Berna e a sua sorella pedofila
Cate, al crucco Valerio e a Flavio il maestro. Vi voglio bene.

\subsection{D}
Ringrazio gli eroi di D. Qui c'è da dire
veramente una gran quantità di cose (perdonatemi!). Questo gruppo fu
originariamente fondato da: me (Gordon Freeman), Andrei (detto Mr Dry,
quello con la dad lore più grande di tutti), Giuseppe (detto er Giuse,
il terrone potentino), Francesco (detto FrancyAI, a cui ci
rivolgiamo quando finiamo i token), Francesco Paoletti (il Francesco
scomparso) e Michele (che ormai fa un sacco di impicci esteri), il primo
giorno dell'Università. Questo gruppo nacque dalla necessità di stretta
sopravvivenza universitaria: bisognava allearsi con i primi con cui si
aveva un minimo di confidenza per scambiarsi informazioni e per avere
qualcuno con cui chiacchierare a lezione. Il primo con cui instaurai
una certa confidenza fu Andrei -- colui che creò effettivamente il
gruppo -- e durante le prime lezioni cazzeggiavamo come pochi: quando
si presentava la prima occasione buona andavamo di corsa in cortile
perché non ne potevamo più di tutti gli scoppiati del nostro corso. Mi
ricordo dei suoi termini strani, delle sue prese in giro per la mia
parlata ternana, delle smagicate, del monopattino in due, della
cioccata di Monti in prima fila perché stava delirando su Brawl Stars.
Da allora fece un grosso glow up, dopo essersi girato mezzo mondo da
solo, collezionando e condividendo un mucchio di esperienze;
sicuramente sarà maturato molto più di me. Ero in buonissimi rapporti
anche con il resto dei membri e già sentivo nell'aria che qualcosa
stava per cambiare.

Con il passare dei mesi il nostro gruppo iniziò ad acquisire forma e
identità. Quella che all'inizio era una lettera messa a caso come
placeholder divenne il nostro simbolo. Si passò a ``D di Domodossola'',
la prima ``season'' canonica in assoluto. Iniziarono pian piano ad
affluire gli altri ``avengers'': Massimo (detto MaxTheOne, il GOAT),
Federico (detto il TheWomen, l'inguaribile laziale), Matteo Greco
(detto Matteo Greek, con quel pizzetto che è diventato suo marchio di
fabbrica), Matteo Lanzillo (detto AURAMAN, il cinofilo per antonomasia,
la cui ossessione batte il talento), Flavio (detto er Bricks, che
adesso sta bruciato come pochi), Mattia (conosciuto come lo chef di
Avezzano, che ci porterà alla laurea in Tailandia), Daniele (detto il
Daniz, anche lui bello GOAT), Vlad (detto boh, non so come chiamatte se
non Vlad, scusa), Francesco (il secondo Francesco attuale, detto il
Moglia, membro onorario del biofregna).

Nel frattempo mi divertivo a cambiare il nome e la foto di gruppo con i
titoli e le immagini più disparate. Tali aggiornamenti scandivano
l'inizio delle ``season'', tutte riguardanti i meme nati durante le
lezioni e le sessioni d'esame; il tutto mantenendo solo un vincolo
fondamentale: D doveva rimanere, come se questa lettera fosse la nostra
stella polare, il carattere dal quale tutto è nato, il nostro big bang,
la nostra inizializzazione in memoria. Seguì poi l'entrata di Lorenzo
(il secondo ternano), che ci mise tanto ad aggiungersi perché
all'inizio mi detestava per via della mia burinagine ternana e quindi
mi teneva a distanza già dal binario della stazione. Il destino volle
che la sua opinione si modificasse e dunque ottenne anche lui il posto
all'interno del gruppo più figo ed esclusivo di tutta la Sapienza (a
mani bassissime). Dopo questa lunga cronaca posso finalmente passare ai
ringraziamenti. Questo gruppo mi ha regalato tante tante gioie; mi ha
svoltato l'intera carriera universitaria: senza, mi sarei laureato
sicuramente almeno per l'anno prossimo, post GTA VI -- una cosa pi\`u che inconcepibile.

La nostra è un'intelligence di cucinati che riesce incredibilmente a
lavorare in armonia: ognuno fa la sua parte e ci diamo forza l'un
l'altro. Mi pento solo di non aver potuto passare più esperienze in
vostra compagnia al di fuori dell'uni, per via dei maledettissimi treni
di Trenitalia. Adesso sei di noi sono arrivati per primi al traguardo
(i GOAT), quelli rimasti hanno ancora un bel po' di gas da dare, ma
sono sicuro che ce la faranno presto. Non ci tireremo mai indietro nel
dare una mano (e i nostri token); siamo uniti fino alla fine e spero
che lo saremo anche oltre, nelle prossime season dei prossimi capitoli.
Vi dico infine con sincerità che vi sarei molto grato se potessimo
votare insieme l'espulsione dei due parassiti rimasti -- di cui non ho
citato il nome --, cosicché la mia anima possa riposare contenta. Un
abbraccio a tutti voi: D di a tutto gas.

\section{AI corporates}
Ultimissimissimissimi ringraziamenti speciali vanno a OpenAI, Google e
Anthropic per avermi fornito i loro potentissimi modelli di
Intelligenza Artificiale, impiegati per la riuscita di questo ultimo
anno e per accelerare significativamente la redazione di questa tesi.

\section{Citazioni}
Ultimissimissimissimissimissimi ringraziamenti ad alcune delle
citazioni -- tratte da libri, film e videogiochi -- che mi hanno
plasmato:

\begin{itemize}

\item ``Considerate la vostra semenza:\\
fatti non foste a viver come bruti,\\
ma per seguir virtute e canoscenza.'' $\sim$ \emph{Dante Alighieri} | Canto XXVI dell'Inferno

\item ``[\dots] che ti fa ciò che quivi si pispiglia?\\
Vien dietro a me, e lascia dir le genti:\\
sta come torre ferma, che non crolla\\
già mai la cima per soffiar di venti.'' $\sim$ \emph{Dante Alighieri} | Canto V
del Purgatorio

\item ``Audentes fortuna iuvat'' $\sim$ \emph{Virgilio} | Libro X dell'Eneide

\item ``È la pazienza, assai più della forza, a fare i grandi
guerrieri.'' $\sim$ \emph{Tenzen} | Le cronache dell'acero e del
ciliegio

\item ``È il tempo che tu hai perduto per la tua rosa che ha fatto la
tua rosa così importante.'' $\sim$ \emph{Volpe} | Piccolo Principe

\item ``Infilarti le piume nel culo non fa di te una gallina!'' $\sim$
\emph{Tyler Durden} | Fight Club

\item ``A volte sono le persone che nessuno immagina possano fare certe
cose, a fare cose che nessuno può immaginare.'' $\sim$ \emph{Alan Turing} |
The Imitation Game

\item ``La vita è come una scatola di cioccolatini\dots non sai mai quale ti capita.'' $\sim$ \emph{Forrest Gump} | Forrest Gump

\item ``Alla fine, cosa distingue uno schiavo da un uomo?
Denaro? Potere? No. Un uomo sceglie, uno schiavo obbedisce'' $\sim$
\emph{Andrew Ryan} | BioShock

\item ``L'uomo giusto nel posto sbagliato può fare una grande
differenza. Perciò si svegli, Mr. Freeman. Apra gli occhi, e si
svegli'' $\sim$ \emph{G-Man} | Half-Life 2

\end{itemize}
